\documentclass[11pt,a4paper]{article}
\usepackage[margin=2.4cm]{geometry}
\usepackage{amsmath,amssymb,amsthm,mathtools}
\usepackage{booktabs}
\usepackage{enumitem}
\usepackage{xcolor}
\usepackage{hyperref}
\usepackage{fontspec}
\usepackage{polyglossia}
\setmainlanguage{russian}
\setotherlanguage{english}
\setmainfont{Times New Roman}
\setmonofont{Consolas}[Scale=0.90]

\hypersetup{
  colorlinks=true,
  linkcolor=blue!70!black,
  citecolor=blue!70!black,
  urlcolor=blue!70!black,
  pdftitle={Проективное исчисление над полем Q(i, sqrt 2): Функциональная полнота Поста, полиномиальная редукция 3-SAT, конструктивные когомологии де Рама и барьеры сложности конечных автоматов на Z_512 в Lean 4},
  pdfauthor={Каратаев Дмитрий Александрович}
}

\theoremstyle{plain}
\newtheorem{theorem}{Теорема}[section]
\newtheorem{lemma}[theorem]{Лемма}
\newtheorem{proposition}[theorem]{Предложение}
\newtheorem{corollary}[theorem]{Следствие}
\theoremstyle{definition}
\newtheorem{definition}[theorem]{Определение}
\newtheorem{remark}[theorem]{Замечание}

\title{\textbf{Функциональная полнота Поста, полиномиальная редукция 3-SAT, конструктивные когомологии де~Рама и барьеры сложности конечных автоматов на $\mathbb{Z}_{512}$ в Lean~4}}
\author{Д.\,А.~Каратаев\thanks{AI-Guiders Research (\url{https://github.com/AI-Guiders}). Формальная верификация в Lean~4 и аналитическая разработка выполнены при участии ИИ-партнёра \emph{Луч} (\emph{Luch}, AI-Guiders).}}
\date{}

\begin{document}
\maketitle

\begin{abstract}
Во второй статье серии устанавливаются булево-логические, теоретико-сложностные, когомологические и теоретико-мерные основания проективного исчисления на фазовом кольце $\mathbb{Z}_{512}$ ($Q_\square(3)=128$, $\Phi_{\mathrm{half}}(3)=256$) и $12$-вершинном комплексе кубооктаэдра $O_2(3)$. Доказано, что на антиподальной подгруппе полюсов $\{0,256\}\subset\mathbb{Z}_{512}$ сложение фаз до оператора проекции (Канал~1, линейный класс Поста~$L$: $\{\mathrm{NOT},\mathrm{XOR}_3\}$) и сложение проекций дискретной экспоненты Эйлера с целочисленным делением на $512$ (Канал~2, нелинейный пороговый класс: $\{\mathrm{OR}_3,\mathrm{AND}_3\}$) образуют функционально полную систему по теореме Поста, сводя выполнимость формулы 3-SAT из $M$ дизъюнктов к условию $\sum_{m=1}^M q_{\mathrm{OR}}(\psi_{m,1},\psi_{m,2},\psi_{m,3}) = 0$. Установлены закон склейки остатка $7\to 1$ в прямой проверке, тождественное зануление линейного градиента Канала~2 на замкнутых $4$-дизъюнктных циклах чётности Канала~1 ($\mathrm{xor4LinearGradientSum}\equiv 0$) и ветвление прообразов $7^T$ вдоль лучей обратного конуса объёма $\mathcal{V}_{O_2}(3,T) = (10T^3+15T^2+11T+3)/3$.

Показано, почему реальная криптографическая стойкость определяется не абстрактной внешней машиной Тьюринга с бесконечной лентой вне системы, а конечными автоматами-наблюдателями, погружёнными \emph{внутрь} самого графа состояний: (1)~любой внутренний решатель ограничен четырьмя формальными барьерами (ёмкостью регистра самовключения Дирихле, законом квадрата слепоты проекции $72^2 = 5\,184$ при переходе от внутреннего $M$-носителя к внешнему $T$-наблюдаемому, задержкой пересечения причинного конуса и обратным возмущением кограницы); (2)~поскольку внешнее $T$-измерение извлекает лишь $1$-битный знак из $72$-битной четырёхканальной ячейки Верле на $\mathbb{Z}_{512}$, оставляя $71$ бит внутренней фазы алгебраически ненаблюдаемым без разрушения квантованной кограницы плакетки $512\mathbb{Z}$, фазовая передача на $M$-носителе реализует встроенный информационно-теоретический блокнот Вернама Шеннона ($H(M\mid T)=H(M)$), стойкость которого не зависит от гипотезы $P\ne NP$. Наконец, доказаны конструктивные тождества когомологий де~Рама--Стокса--Гаусса--Черна ($d_1\circ d_0\equiv 0$, $d_2\circ d_1\equiv 0$), обращение нильпотентных матриц без деления ($N_{\mathrm{nil}}^3=0$) и конечно-аддитивная мера без парадокса Банаха--Тарского в ядре \textbf{Lean~4} (\texttt{thm4}, \texttt{thm5}, \texttt{0 sorryAx}, \texttt{0 Classical.choice}).

\medskip\noindent
\textbf{Ключевые слова:} функциональная полнота Поста, редукция 3-SAT, барьеры внутреннего наблюдателя, шифр Вернама Шеннона, дискретные когомологии де~Рама, нильпотентные струи, конструктивная мера, Lean~4.
\end{abstract}

\section{Функционально полный булев базис Поста, представление 3-SAT на треугольных гранях \texorpdfstring{$\mathbb{Z}_{512}$}{Z512} и асимметрия проверки и обращения}
\label{sec:boolean-3sat}

На полюсах $\{0,256\}\subset\mathbb{Z}_{512}$ ($\mathcal{C}_3(0)=+128$, $\mathcal{C}_3(256)=-128$) Канал~1 задаёт $\mathrm{NOT}(\psi)=(\psi+256)\bmod 512$ и чётность $\mathrm{XOR}_3$ по закону $Q_\square(d)^2\mathcal{C}_d(\psi_1+\psi_2+\psi_3)\equiv \mathcal{C}_d(\psi_1)\mathcal{C}_d(\psi_2)\mathcal{C}_d(\psi_3)$, а Канал~2 через дефект $\Delta_{\mathrm{3SAT}}(\psi_1,\psi_2,\psi_3) := Q_\square(d)-(\mathcal{C}_d(\psi_1)+\mathcal{C}_d(\psi_2)+\mathcal{C}_d(\psi_3))$ задаёт $q_{\mathrm{OR}} := \mathrm{tdiv}(\Delta_{\mathrm{3SAT}},512)\in\{0,1\}$ и $q_{\mathrm{AND}} := \mathrm{tdiv}(Q_\square(d)+\sum_{k=1}^3\mathcal{C}_d(\psi_k),512)\in\{0,1\}$.

\begin{theorem}[Функциональная полнота Поста, представление 3-SAT и склейка остатка $7\to 1$]
\label{thm:main-4}
Для всех $d\in\mathbb{N}$ и всех $p_1,p_2,p_3\in\mathbb{Z}$ дефект $\Delta_{\mathrm{3SAT}} = Q_\square(d) - (\mathcal{C}_d(p_1)+\mathcal{C}_d(p_2)+\mathcal{C}_d(p_3))$ тождественно раскладывается в виде (\texttt{thm4\_boolean\_post\_completeness\_3sat\_reduction\_and\_channel\_orthogonality}):
\[
  \Delta_{\mathrm{3SAT}}(d;p_1,p_2,p_3) \;=\; N_{\mathrm{ring}}(d)\,q_{\mathrm{OR}}(d;p_1,p_2,p_3) \;+\; r_{\mathrm{drop}}(d;p_1,p_2,p_3).
\]
При $d=3$ на полюсах $\{0,256\}^3\subset\mathbb{Z}_{512}^3$:
\begin{enumerate}[label=\textbf{(\roman*)},leftmargin=*,noitemsep]
\item частное $q_{\mathrm{AND}}$ равно~$1$ исключительно на $(0,0,0)$, а частное $q_{\mathrm{OR}}$ равно~$0$ на всех $7$~выполняющих октантах дизъюнкта $(l_1\lor l_2\lor l_3)$ и равно~$1$ исключительно на $(256,256,256)$, так что выполнимость формулы 3-SAT из $M$ дизъюнктов эквивалентна $\sum_{m=1}^M q_{\mathrm{OR}}(\psi_{m,1},\psi_{m,2},\psi_{m,3}) = 0$;
\item прямое вычисление $q_{\mathrm{OR}}$ отбрасывает остаток $r_{\mathrm{drop}}\in\{-256,0,+256\}$ и склеивает $2^3-1=7$ выполняющих октантов в один выход $q_{\mathrm{OR}}=0$ (ветвление прообраза $1\to 7$ на каждой грани).
\end{enumerate}
\end{theorem}

\begin{theorem}[Некоммутативность каналов, кубическая рекуррентность обратного конуса и слепая многоузловая телепортация]
\label{thm:main-4b}
Для всех $T,\psi,a,r_{1,\mathrm{out}},r_{2,\mathrm{out}},P\in\mathbb{Z}$ (\texttt{thm4\_boolean\_post\_completeness\_3sat\_reduction\_and\_channel\_orthogonality}):
\begin{enumerate}[label=\textbf{(\roman*)},leftmargin=*,noitemsep]
\item для четвёрки дизъюнктов 3-SAT, кодирующей линейное уравнение чётности Канала~1 ($x_1\oplus x_2\oplus x_3=0$), сумма линейных проекций $\mathcal{C}_3$ тождественно равна нулю ($\mathrm{xor4LinearGradientSum}\equiv 0$), а операторы $\mathcal{C}_3$ и $q_{\mathrm{AND}}$ недистрибутивны относительно сложения фаз Канала~1;
\item при добавлении $(T+1)$-й оболочки кубооктаэдра $10(T+1)^2+2$ к утроенному объёму шара $3\mathcal{V}_{O_2}(3,T) = 10T^3+15T^2+11T+3$ выполняется тождество
  \[
    (10T^3+15T^2+11T+3) + 3\bigl(10(T+1)^2+2\bigr) \;=\; 10(T+1)^3 + 15(T+1)^2 + 11(T+1) + 3,
  \]
  с ветвлением лучевых прообразов $7^3 = 343 > 2^{6-3} = 8$ при $T=3$;
\item для антиподальных пар $r_{1,\mathrm{in}} = a+P$, $r_{2,\mathrm{in}} = r_{1,\mathrm{out}}+P$, $b = r_{2,\mathrm{out}}+P$ и слепых синдромов $\sigma_A = \psi-a-P$, $\sigma_{R_1} = r_{1,\mathrm{in}}-r_{1,\mathrm{out}}-P$, $\sigma_{R_2} = r_{2,\mathrm{in}}-r_{2,\mathrm{out}}-P$ в сети $A\leftrightarrow R_1\leftrightarrow R_2\leftrightarrow B$ сумма на приёмнике тождественно сокращает все промежуточные ключи: $b + \sigma_{R_2} + \sigma_{R_1} + \sigma_A = \psi$.
\end{enumerate}
\end{theorem}

\section{Внутренние конечные вычислители против внешних машин Тьюринга и встроенный блокнот Вернама Шеннона на \texorpdfstring{$\mathbb{Z}_{512}$}{Z512}}
\label{sec:embedded-adversary-vernam}

В терминологии классической теории моделей А.\,И.~Мальцева (\emph{«Алгебраические системы»}~\cite{Malcev1970}), где \emph{носителем} (\emph{carrier}) называется основное множество алгебраического сорта, рассматриваемое исчисление образует двусортную алгебраическую систему $\langle\mathcal{M},\mathcal{T};\,\pi_{M\to T},\,\iota_{T\to M}\rangle$ с \textbf{дискретным $M$-носителем} $\mathcal{M}=\mathbb{Z}_{512}$ (внутренние модулярные фазовые состояния на решётке кубооктаэдра) и \textbf{непрерывным $T$-носителем} $\mathcal{T}$ (внешне наблюдаемые пороговые и евклидовы инварианты).

В классической теории сложности взломщик моделируется абстрактной машиной Тьюринга, расположенной \emph{снаружи} фазового пространства с бесконечной внешней лентой и бесплатным чтением состояния. Из Теорем~\ref{thm:main-4} и~\ref{thm:main-5} следуют два структурных результата для любого \emph{реализуемого} вычислителя:
\begin{enumerate}[label=\textbf{(\arabic*)},leftmargin=*]
\item \textbf{Любой реализуемый вычислитель является конечным подавтоматом внутри носителя и подчинён четырём формальным барьерам обращения:}
  конечный подкомплекс из $V_{\mathrm{adv}} < \mathcal{V}_{O_2}(3,T)$ ячеек внутри причинного конуса ограничен:
  (I)~барьером Дирихле самовключения регистра ($B_2 V_{\mathrm{adv}} < B_2\mathcal{V}_{O_2}(3,T)$ при $B_2 = 8\times 9 = 72$ бит на ячейку);
  (II)~законом квадрата слепоты проекции $K_P(3)\cdot 81 = 72^2 = 5\,184$ (грубый замер на $T$-слое извлекает лишь $1$~бит знака из $72$ внутренних бит ячейки $M$-носителя);
  (III)~задержкой распространения по конусу ($2T > T$);
  (IV)~обратным возмущением кограницы $(d_1 r)_f = 512\,n_f$ при записи в регистр.
\item \textbf{Встроенный информационно-теоретический блокнот Вернама на дискретном $M$-носителе $\mathbb{Z}_{512}$:}
  поскольку внешнее наблюдение на $T$-слое видит лишь $1$-битное пороговое частное, тогда как остальные $71$~бит внутренней фазы ячейки (и остаток $r\in\{-256,0,+256\}$ на ядрах чётности с $\mathrm{xor4LinearGradientSum}\equiv 0$) остаются алгебраически ортогональными $T$-проекции и топологически защищёнными целочисленным классом Черна $(d_1 r)_f = 512\,n_f$, фазовая передача на $M$-носителе обладает встроенным режимом одноразового блокнота Вернама ($H(M\mid T) = H(M)$). По теореме Шеннона совершенная секретность блокнота Вернама не может быть нарушена никаким вычислителем --- независимо от разрешения гипотезы $P$ vs.\ $NP$ для абстрактных внешних лент Тьюринга.
\end{enumerate}

\section{Конструктивные когомологии де~Рама--Стокса--Гаусса--Черна, линейная алгебра без деления и мера без парадокса Банаха--Тарского}
\label{sec:constructive-foundations}

\begin{theorem}[Евклидово расщепление связности и конструктивные когомологии де~Рама--Стокса--Гаусса--Черна]
\label{thm:main-5}
Для всех $d\in\mathbb{N}$, всех разностей фаз $\Delta\varphi\in\mathbb{Z}$, всех $0$-коцепей $(p_0,p_1,p_2,p_3)\in\mathbb{Z}^4$ и всех $1$-коцепей $(e_1,\dots,e_{12})\in\mathbb{Z}^{12}$ на $12$~рёбрах куба (\texttt{thm5\_constructive\_de\_rham\_cohomology\_nilpotent\_algebra\_and\_banach\_tarski\_free\_measure}):
\begin{enumerate}[label=\textbf{(\roman*)},leftmargin=*,noitemsep]
\item частное $q = \mathrm{branchQuot}(d,\Delta\varphi)$ и остаток $r = \mathrm{edgeConn}(d,\Delta\varphi)$ удовлетворяют $N_{\mathrm{ring}}(d)\,q + r \equiv \Delta\varphi$;
\item на любом замкнутом $4$-цикле грани~$f$ кограница $(d_1(d_0 p))_f = (p_1-p_0)+(p_2-p_1)+(p_3-p_2)+(p_0-p_3) \equiv 0$, а циркуляция связности квантована периодом кольца: $(d_1 r)_f = N_{\mathrm{ring}}(d)\,n_f \in N_{\mathrm{ring}}(d)\mathbb{Z}$ ($n_f := -\sum_{e\in\partial f} q_e\in\mathbb{Z}$);
\item сумма потоков $(d_1 e)_f$ по всем $6$~граням замкнутой поверхности куба тождественно равна нулю ($d_2\circ d_1\equiv 0$).
\end{enumerate}
\end{theorem}

\begin{theorem}[Обращение нильпотентных и $\mathrm{SL}_2(\mathbb{Z})$ симплектических операторов без деления]
\label{thm:main-5b}
Для всех $(s_0,s_1,s_2,s_3)\in\mathbb{Z}^4$ и всех $(u,v,\varphi)\in\mathbb{Z}^3$ (\texttt{thm5\_constructive\_de\_rham\_cohomology\_nilpotent\_algebra\_and\_banach\_tarski\_free\_measure}):
\begin{enumerate}[label=\textbf{(\roman*)},leftmargin=*,noitemsep]
\item 4D нильпотентный оператор $2\mathcal{M}(s_0,s_1,s_2,s_3) = (2s_0+s_1,\,2s_1+2s_2,\,2s_2,\,2s_3)$ ($N_{\mathrm{nil}}^3 = 0$, $\det(2\mathcal{M})=24$) точно обращается трёхчленным полиномом Неймана $12(I_4 - N_{\mathrm{nil}} + N_{\mathrm{nil}}^2)(2\mathcal{M}\,s) = (24s_0,24s_1,24s_2,24s_3)$;
\item симплектический шаг Верле $(u',v') = (2u-v+\varphi,\,u)$ ($\det=+1\in\mathrm{SL}_2(\mathbb{Z})$) точно обращается формулой $(v',\,2v'-u'+\varphi) = (u,v)$.
\end{enumerate}
\end{theorem}

\begin{theorem}[Вариационная оптимизация на $L^1$-грани, аддитивность меры без аксиомы выбора и законы цепей переменного/постоянного тока]
\label{thm:main-5c}
Для всех $Q,k,\mathcal{C}_\varphi,\mathcal{S}_\varphi\in\mathbb{Z}$ (\texttt{thm5\_constructive\_de\_rham\_cohomology\_nilpotent\_algebra\_and\_banach\_tarski\_free\_measure}):
\begin{enumerate}[label=\textbf{(\roman*)},leftmargin=*,noitemsep]
\item сопряжённые проекции $V = 2Q-k$ и $I = k$ удовлетворяют $V+I \equiv 2Q$, а квадратичная форма $P = Q^2 - (k-Q)^2 \le Q^2$ достигает максимума $Q^2$ при $k^* = Q$;
\item послойное разбиение шара кубооктаэдра $\mathcal{B}_{O_2}(3,3)$ сохраняет целочисленную меру $1 + 12 + 42 + 92 = 147$ без аксиомы выбора (\texttt{0 Classical.choice}), а на кольце $\mathbb{Z}_{512}$ реактивная мощность равна нулю ($\sum\mathcal{C}_3\mathcal{S}_3 = 0$), однополупериодное выпрямление даёт $Q_\square(3)^2 = 16\,384$, двухполупериодный мост Греца даёт $2Q_\square(3)^2 = 32\,768$, и противофазные токи LC-резонанса тождественно компенсируются ($-\mathcal{S}_\varphi + \mathcal{S}_\varphi = 0$, $\mathcal{C}_\varphi + (-\mathcal{C}_\varphi) = 0$).
\end{enumerate}
\end{theorem}

\section{Формальная верификация в Lean~4}
Формальное доказательство Теорем~\ref{thm:main-4} и~\ref{thm:main-5} (\texttt{thm4}, \texttt{thm5}) содержится в модуле \texttt{LogicCohomology.lean} в репозитории \url{https://github.com/AI-Guiders/scientific-papers} (\texttt{0 sorryAx}, \texttt{0 Classical.choice}).

\begin{thebibliography}{9}
\bibitem{Malcev1970}
А.\,И.~Мальцев, \emph{Алгебраические системы}, М.: Наука, 1970 (англ.\ пер.: A.\,I.~Mal'cev, \emph{Algebraic Systems}, Springer-Verlag, 1973).
\bibitem{Post1941}
E.\,L.~Post, \emph{The Two-Valued Iterative Systems of Mathematical Logic}, Princeton University Press, 1941.
\bibitem{Shannon1949}
C.\,E.~Shannon, ``Communication theory of secrecy systems,'' \emph{Bell System Technical Journal} \textbf{28}(4) (1949), 656--715.
\bibitem{Moura2021}
L.~de~Moura and S.~Ullrich, ``The Lean~4 theorem prover and programming language,'' in \emph{CADE~28}, LNCS vol.~12699, Springer, 2021, pp.~625--635.
\end{thebibliography}

\end{document}
